\section{Momentos genealógicos, operadores de Jacobi y residuos espectrales}
\label{sec:jacobi-generated}

La medida inducida por el calendario proporciona un enlace entre la
geometría de las celdas y una realización espectral positiva. Este enlace
utiliza toda la partición marcada: cada uno de sus doce intervalos conserva
la misma masa cronológica, mientras su longitud procede de la coordenada
correspondiente del registro. La densidad compensa exactamente esa
longitud. Por consiguiente, la distribución espacial de la medida retiene
la anisotropía del registro aunque la distribución entre ventanas sea
uniforme.

Sean \(K_j>0\), \(Q=\sum_{j=1}^{12}K_j\), \(d_j=K_j/Q\) y
\(t_j=\sum_{i\leq j}d_i\), con \(t_0=0\). La medida de probabilidad
utilizada en esta sección es
\[
 \dd\mu_K(t)=\sum_{j=1}^{12}
 \frac{\mathbf1_{[t_{j-1},t_j]}(t)}{12d_j}\,\dd t.
\]
Los extremos compartidos tienen medida nula. Sus momentos son racionales
para todo registro entero:
\begin{equation}
 m_r=\int_0^1t^r\,\dd\mu_K(t)
 =\frac1{12(r+1)}\sum_{j=1}^{12}
   \frac{t_j^{r+1}-t_{j-1}^{r+1}}{d_j},\qquad r\geq0.
 \label{eq:jacobi-moments}
\end{equation}
La igualdad \(m_0=1\) expresa conservación de la masa total. Para cada
\(n\geq1\), la matriz de Hankel \(H_n=(m_{a+b})_{0\leq a,b<n}\)
es definida positiva: su forma cuadrática es la integral del cuadrado
de un polinomio no nulo sobre un intervalo donde la densidad es positiva.
Esta observación permite pasar de los momentos a un operador sin elegir
un espectro objetivo.

\begin{theorem}[Recurrencia tridiagonal determinada por la partición]
\label{thm:jacobi-recurrence}
Existe una única sucesión de polinomios mónicos \(p_n\), de grado \(n\),
ortogonales para \(\mu_K\). Si
\[
 h_n=\int p_n^2\,\dd\mu_K>0,\qquad
 a_n=\frac{\int t p_n^2\,\dd\mu_K}{h_n},\qquad
 \beta_n=\frac{h_n}{h_{n-1}}>0\quad(n\geq1),
\]
entonces \(p_0=1\), \(p_1=t-a_0\) y
\begin{equation}
 p_{n+1}(t)=(t-a_n)p_n(t)-\beta_np_{n-1}(t),\qquad n\geq1.
 \label{eq:jacobi-recurrence}
\end{equation}
Todos los coeficientes \(a_n,\beta_n\) son racionales para el registro
entero declarado. En la base ortonormal \(\psi_n=p_n/\sqrt{h_n}\),
la compresión de la multiplicación por \(t\) a los polinomios de grado
menor que \(n\) tiene matriz simétrica
\[
 J_n=\begin{pmatrix}
 a_0&\sqrt{\beta_1}&&0\\
 \sqrt{\beta_1}&a_1&\ddots&\\
 &\ddots&\ddots&\sqrt{\beta_{n-1}}\\
 0&&\sqrt{\beta_{n-1}}&a_{n-1}
 \end{pmatrix}.
\]
\end{theorem}
\begin{proof}
La positividad de \(H_n\) garantiza que el sistema lineal que impone
ortogonalidad a un polinomio mónico tiene solución única. Como sus
coeficientes son momentos racionales, la solución es racional. Se
descompone \(tp_n\) en la base mónica hasta grado \(n+1\). Para
\(j\leq n-2\), la simetría del producto interno da
\(\langle tp_n,p_j\rangle=\langle p_n,tp_j\rangle=0\).
Quedan únicamente los términos de grados \(n+1,n,n-1\). Sus
coeficientes son, respectivamente, \(1,a_n,h_n/h_{n-1}\), lo que
demuestra la recurrencia. La normalización por \(\sqrt{h_n}\)
convierte los dos coeficientes adyacentes en \(\sqrt{\beta_n}\).
\end{proof}

Para el registro expresado en la sección~\ref{sec:register}, la primera
posición diagonal y el cuadrado del primer acoplamiento son
\[
 a_0=m_1=\frac{71195}{150312},\qquad
 \beta_1=m_2-m_1^2=\frac{2206226167}{22593697344}.
\]
La primera igualdad calcula el centro de masa de la partición cronológica;
la segunda calcula su dispersión y determina el primer enlace de la
matriz tridiagonal. Ambas cifras se obtienen integrando las longitudes
generadas, y su significado estructural antecede a cualquier evaluación
decimal.

\begin{theorem}[Espectro simple, cuadratura y residuos positivos]
\label{thm:jacobi-residues}
Los autovalores \(\lambda_1<\cdots<\lambda_n\) de \(J_n\) son
simples y pertenecen a \((0,1)\). Existen pesos únicos
\(w_i>0\), de suma uno, tales que
\begin{equation}
 \int f(t)\,\dd\mu_K(t)=\sum_{i=1}^{n}w_i f(\lambda_i)
 \qquad(\deg f\leq2n-1).
 \label{eq:jacobi-quadrature}
\end{equation}
El resolvente observado en la constante unitaria \(e_0\) satisface
\begin{equation}
 g_n(z)=\langle e_0,(zI-J_n)^{-1}e_0\rangle
       =\sum_{i=1}^{n}\frac{w_i}{z-\lambda_i}
       =\cfrac1{z-a_0-\cfrac{\beta_1}{z-a_1-
          \cfrac{\beta_2}{\ddots-\cfrac{\beta_{n-1}}{z-a_{n-1}}}}}.
 \label{eq:jacobi-resolvent}
\end{equation}
En particular, sus polos y residuos se obtienen del registro por las
operaciones de momento, ortogonalización y compresión.
\end{theorem}
\begin{proof}
Para cualquier polinomio no nulo \(p\) de grado menor que \(n\),
\[
 0<\int_0^1 t|p(t)|^2\,\dd\mu_K(t)
   <\int_0^1 |p(t)|^2\,\dd\mu_K(t).
\]

La caracterización variacional coloca el espectro de \(J_n\) en
\((0,1)\). Cada subdiagonal es positiva. Las ecuaciones de un autovector
determinan recursivamente todas sus coordenadas a partir de la primera;
si ésta se anulase, se anularía el vector entero. Cada autoespacio tiene
así dimensión uno. La simetría de la matriz da diagonalización ortogonal
y simplicidad del espectro. Para autovectores unitarios \(v_i\), se toma
\(w_i=|\langle e_0,v_i\rangle|^2>0\). La completitud de la base
propia da \(\sum_iw_i=1\) y la primera expresión del resolvente.

La recurrencia tridiagonal muestra por expansión del determinante que
\(\det(tI-J_n)=p_n(t)\). Para \(\deg f\leq2n-1\), la división
euclídea escribe \(f=qp_n+r\), con \(\deg q,\deg r<n\).
La ortogonalidad anula la integral de \(qp_n\), y su evaluación en los
\(\lambda_i\) también se anula. Para \(r\) de grado menor que
\(n\), las sucesivas multiplicaciones de la constante por \(t\)
permanecen dentro del espacio de polinomios hasta el grado requerido;
por ello \(\int r\,\dd\mu_K=\langle e_0,r(J_n)e_0\rangle\).
La resolución espectral prueba la fórmula de cuadratura. Su unicidad
se deduce de la invertibilidad de la matriz de Vandermonde evaluada
en las raíces distintas. Finalmente, el complemento de Schur de la
primera fila y columna de \(zI-J_n\), iterado sobre las submatrices
terminales, produce la fracción continua.
\end{proof}

\begin{figure}[htbp]
\centering
\begin{tikzpicture}[>=Latex, node distance=8mm and 8mm,
 every node/.style={font=\small}, box/.style={draw,rounded corners=1pt,
 text width=37mm,minimum height=13mm,align=center}]
\node[box] (k) {Registro marcado\\\(K,Q\)};
\node[box,right=of k] (m) {Medida cronológica\\\(\mu_K\)};
\node[box,right=of m] (h) {Momentos positivos\\\(H_n\)};
\node[box,below=of h] (j) {Operador tridiagonal\\\(J_n\)};
\node[box,left=of j] (g) {Resolvente\\\(g_n(z)\)};
\node[box,left=of g] (w) {Polos y residuos\\\(\lambda_i,w_i>0\)};
\draw[->] (k)--(m);\draw[->] (m)--(h);\draw[->] (h)--(j);
\draw[->] (j)--(g);\draw[->] (g)--(w);
\end{tikzpicture}
\caption{Realización espectral de la partición HMT. La geometría de las
celdas determina la medida; la positividad de sus momentos produce el
operador y los residuos positivos de su resolvente.}
\label{fig:jacobi-generated}
\end{figure}

La fórmula del resolvente aporta una lectura espectral de la memoria
espacial. Los pesos positivos describen cuánto observa la coordenada
constante de cada modo; los polos son las frecuencias algebraicas del
operador de multiplicación comprimido. Los momentos de orden creciente
recuperan detalles sucesivos de la medida, y los coeficientes de la
fracción continua conservan esa información mediante una recurrencia
local de tres términos. La apariencia clásica de una matriz de Jacobi
se reconoce después de haber construido sus coeficientes desde las
celdas HMT. La cadena causal queda fijada por
\[
 K\longmapsto(d,t)\longmapsto\mu_K\longmapsto(m_r)_{r\geq0}
 \longmapsto\bigl((a_n)_{n\geq0},(\beta_n)_{n\geq1}\bigr)
 \longmapsto(J_n,g_n).
\]

\begin{proposition}[Compatibilidad de dimensión y de refinamiento]
\label{prop:jacobi-refinement}
Las matrices \(J_n\) son compresiones principales compatibles de la
misma multiplicación acotada por \(t\) en \(L^2(\mu_K)\). Sus
medidas espectrales \(\sum_iw_i\delta_{\lambda_i}\) convergen
débilmente a \(\mu_K\). Si cada subintervalo nonádico se representa
por su punto medio y por su masa \(1/(12\cdot9^r)\), la medida
atómica resultante \(\mu_{K,r}\) aproxima los momentos de
\(\mu_K\). Para un grado fijo \(n\), la construcción desde esos
momentos recupera los coeficientes \(a_j,\beta_j\) de la misma
matriz cuando el refinamiento tiende a profundidad infinita.
\end{proposition}
\begin{proof}
La ortogonalización de la sucesión completa utiliza las mismas fórmulas
al aumentar \(n\), de modo que cada bloque principal permanece fijo.
La cuadratura reproduce exactamente cualquier polinomio cuando
\(2n-1\) supera su grado. La aproximación uniforme de funciones
continuas en \([0,1]\) por polinomios, junto con la masa total uno,
prueba la convergencia débil. Las sumas de puntos medios convergen
en cada uno de los doce intervalos con su densidad constante, y por
tanto aproximan todos los momentos. En dimensión fija, los coeficientes de
ortogonalización son funciones racionales de un número finito de
momentos, con denominadores formados por determinantes positivos de
Gram. La convergencia de esos momentos y la positividad del límite
aseguran la convergencia de los coeficientes.
\end{proof}

\begin{falsifier}
Una longitud nula invalida la densidad declarada y puede destruir el
rango de los momentos. La sustitución de la medida fija por conteos
no normalizados en mallas distintas cambia el operador. Los residuos
de \(g_n\) son pesos espectrales del resolvente, mientras los residuos
de una forma canónica se definen sobre divisores geométricos: una
identificación entre ambas nociones exige especificar la forma y el
mapa correspondiente. Del mismo modo, \(J_n\) es adimensional; su
identificación con un Hamiltoniano físico requiere la representación
y la escala de acción y reloj que utilice ese sector.
\end{falsifier}

\begin{theorem}[Fidelidad de la realización espectral marcada]
\label{thm:jacobi-marked-faithfulness}
El operador infinito de Jacobi \(J\), junto con su vector cíclico
unitario \(e_0\) y la carga \(Q\), determina el registro \(K\).
Dos realizaciones de este tipo son unitariamente equivalentes mediante
una aplicación que conserva el vector marcado y la carga si y sólo
si sus registros coinciden. El espectro como conjunto es, en cambio,
\([0,1]\) para todo registro positivo del dominio.
\end{theorem}
\begin{proof}
Los polinomios son densos en \(L^2(\mu_K)\): primero se aproximan
funciones continuas y después se usa su densidad para la medida finita.
La base ortonormal construida define una aplicación unitaria que lleva
\(e_n\) a \(\psi_n\). Bajo ella, \(J\) es la multiplicación por
\(t\), un operador autoadjunto acotado, y \(e_0\) representa la
función constante uno. Por tanto,
\[
 \langle e_0,J^re_0\rangle=m_r\qquad(r\geq0).
\]
Estos momentos determinan la medida: dos medidas con los mismos
momentos coinciden en polinomios, después en funciones continuas por
aproximación uniforme y finalmente como medidas de Borel.
Su distribución \(F(t)=\mu_K([0,t])\) es continua y estrictamente
creciente. Cada intervalo inicial tiene masa \(1/12\), de modo que
\[
 t_j=F^{-1}(j/12),\qquad
 K_j=Q\bigl(F^{-1}(j/12)-F^{-1}((j-1)/12)\bigr).
\]
Una equivalencia unitaria que conserva \(e_0\) conserva los momentos
y por ello el registro. La recíproca se sigue de la construcción única.

Para el último enunciado, la densidad de \(\mu_K\) es positiva
en todo subintervalo. Cada \(\lambda\in[0,1]\) admite funciones
unitarias concentradas en intervalos cada vez menores alrededor de
\(\lambda\), sobre las que la norma de \((J-\lambda I)f\)
tiende a cero. Así todo el intervalo pertenece al espectro. Fuera de
él, la multiplicación por \(1/(t-\lambda)\) es acotada y define
el resolvente. Esto prueba la igualdad del espectro.
\end{proof}

El teorema distingue la extensión del espectro y la información que
éste conserva respecto de una observación marcada. El mismo intervalo
espectral admite las distintas distribuciones de la colección; la medida
observada desde la unidad recupera las separaciones originales. El
enlace con la representación grassmanniana es ahora reversible:
los menores marcados reconstruyen \(K\), el registro construye
\((J,e_0,Q)\), y sus momentos y cuantiles recuperan el mismo registro.
Cada lector que factoriza por \((K,\xi)\) posee, por consiguiente,
una expresión espectral marcada una vez conservadas las restantes
coordenadas \(\xi\). La realización espectral contiene estructura,
además de posiciones de polos o valores de autovalores.
